% =====================================================================
% CAPITULO 8 --- WRAPPER VISUAL DIVERSIFICADO
% Mantem a teoria do capitulo e substitui o banco complementar pelo banco
% diversificado, com circuitos e valores proprios em cada questao pratica.
% =====================================================================
\begingroup

% Questões do corpo do capítulo que dependem explicitamente de equivalentes
% calculados anteriormente repetem apenas o diagrama necessário.
\def\fvThCapCargaSeis{%
\begin{circuitikz}[american,scale=.88,transform shape,line width=.8pt]
  \draw (0,0) to[\fonteV,l^={$V_{Th}=\SI{6}{\volt}$}] (0,3.6)
        to[R,l={$R_{Th}=\SI{2}{\ohm}$}] (4.0,3.6)
        to[R,l={$R_L=\SI{6}{\ohm}$},i>^={$I_L$}] (4.0,0) -- (0,0);
\end{circuitikz}}

\def\fvThCapConversao{%
\begin{circuitikz}[american,scale=.72,transform shape,line width=.8pt]
  \draw (0,0) to[\fonteV,l^={$\SI{6}{\volt}$}] (0,3.3)
        to[R,l={$\SI{2}{\ohm}$}] (3.2,3.3) -- (4.1,3.3) node[ocirc]{};
  \draw (0,0)--(4.1,0) node[ocirc]{};
  \node[font=\bfseries] at (2.0,4.0) {Thévenin};
  \draw (7.0,0) to[isource,l^={$\SI{3}{\ampere}$}] (7.0,3.3) -- (10.2,3.3) -- (11.1,3.3) node[ocirc]{};
  \draw (10.2,3.3) to[R,l_={$\SI{2}{\ohm}$}] (10.2,0);
  \draw (7.0,0)--(11.1,0) node[ocirc]{};
  \node[font=\bfseries] at (9.0,4.0) {Norton};
\end{circuitikz}}

\def\fvThCapCargasMultiplas{%
\begin{circuitikz}[american,scale=.86,transform shape,line width=.8pt]
  \draw (0,0) to[\fonteV,l^={$V_{Th}=\SI{18}{\volt}$}] (0,3.7)
        to[R,l={$R_{Th}=\SI{3}{\ohm}$}] (4.0,3.7)
        to[R,l={$R_L$},i>^={$I_L$}] (4.0,0) -- (0,0);
  \node[right,align=left,font=\footnotesize] at (5.0,2.0)
    {$R_L=\SI{3}{\ohm}$\\$R_L=\SI{6}{\ohm}$\\$R_L=\SI{9}{\ohm}$};
\end{circuitikz}}

\newcount\fvThCapQuestao
\fvThCapQuestao=0
\let\fvThCapRespostaOriginal\resposta
\renewcommand{\resposta}[1]{%
  \global\advance\fvThCapQuestao by 1\relax
  \ifnum\fvThCapQuestao=4\relax\circuitoq{\fvThCapCargaSeis}\fi
  \ifnum\fvThCapQuestao=5\relax\circuitoq{\fvThCapConversao}\fi
  \ifnum\fvThCapQuestao=8\relax\circuitoq{\fvThCapCargasMultiplas}\fi
  \fvThCapRespostaOriginal{#1}}

\let\fvThInputOriginal\input
\renewcommand{\input}[1]{%
  \ifstrequal{#1}{bancos/banco-teoremas-de-thevenin-e-norton.tex}%
    {\begingroup\let\input\fvThInputOriginal
     \fvThInputOriginal{bancos/banco-teoremas-de-thevenin-e-norton-diverso.tex}%
     \endgroup}%
    {\fvThInputOriginal{#1}}}
\fvThInputOriginal{cap08-thevenin-norton}
\endgroup
